\pdfbookmark[1]{On Pronunciation}{IPA}

\subsection*{On Pronunciation}

There are different dialects spoken on Pandora (usually referred to by clan or area where it is spoken). It is needless to say that one is not better than the other or that there is anything like a standard or preferred dialect. 

Since the first contact between humans and Na'vi presumably happened with the Omatikaya Clan, the study of the language also started with the Omatikayan Dialect, which is usually referred to as `Forest Na'vi' (\emph{Lì'fya Na'rìngä}). Its pronunciation and vocabulary is the best known variety of Na'vi (so far) and is therefore predominantly concentrated on in this dictionary. Every pronunciation refers to Forest Na'vi first, possible variations are noted after the IPA or directly in the lexical entry.

\subsection*{The Sounds of Na'vi -- IPA}

The Na'vi alphabet has \on{33} distinct sounds that are written with the following letters:
\begin{center}
\B{'  a  aw  ay  ä  e  ew  ey  f  h  i  ì  k  kx  l  ll m \\ 
n  ng  o  p  px  r  rr  s  t  tx  ts  u  v  w  y  z}
\end{center}

The letters b, c, d, j, and q don't exist, g and x only appear in combination---they don't have a sound value of their own.

The pronunciaton of Na'vi is quite regular, although for an English speaker, it has some unusual sounds and sound combinations that can best be learned by listening to sound samples.

\subsubsection*{Vowels}

\begin{tabular}{ c l c }
    &  & IPA \\
  \B{a} & as in `father' & \I{a} \\
  \B{ä} & as in `cat' & \I{\ae} \\
  \B{e} & as in `when' & \I{E} \\
  \B{i} & as the long /i/ in `feat' or `beat' & \I{i} \\
  \B{ì} & as the short /i/ in `fit' or `bit' & \I{I} \\
  \B{o} & as in `so', `or' and `fort' & \I{o} \\
  \B{u} & as in `food' and `boot' (in open and closed syllables) & \I{u} \\ 
    & or `foot' and `book' (in closed syllables) & \I{U} \\
\end{tabular} \\

Note: The difference between open and closed syllables is whether the syllable ends in a vowel (open), e.g. \B{lu} \I{[lu]} `be' or a consonant (closed), e.g. \B{tsun} \I{[\t{ts}un]} or \I{\t{[ts}Un]} `can, be able to'.

\subsubsection*{Diphthongs}

\begin{tabular}{ r l l l r l l}
    &  & IPA  \\
  \B{aw} & as in `house' & \I{aw} \\
  \B{ay} & as in `my' & \I{aj} \\
  \B{ew} & no exact English equivalent; close to the /eow/ in `Beowulf' & \I{Ew} \\
  \B{ey} & as in `hey' or `say' & \I{Ej} \\
\end{tabular}

\subsubsection*{Pseudovowels}

\begin{tabular}{ c l c }
    &  & IPA \\
  \B{ll} & light l as in `fill' or `people' & \I{\s{l}} \\
  \B{rr} & very strongly trilled at the hard gums behind the front teeth & \I{\s{r}} \\
\end{tabular} \\

Note: In contrast to the vowels, \B{ll} and \B{rr} can never start a syllable. They must be preceded by a consonant, e.g. \B{'rr}.pxom, sä.\B{pll}.txe or consonant cluster, e.g. nì.\B{fkrr}.

\subsubsection*{Consonants}

\begin{tabular}{ r l l }
    &  & IPA \\
  \B{'} & = `glottal stop', the sharp break in `uh-oh' & \I{P} \\
  \B{f} & as in `fill' or `cliff' & \I{f} \\
  \B{h} & never silent, always pronounced & \I{h} \\
  \B{k} & unaspirated as the k in `skin' at the start of a syllable & \I{k} \\
             & unreleased at the end of a syllable & \I{k\textcorner} \\
  \B{kx} & pronounced k while holding breath & \I{k'} \\
  \B{l} & as in `life' or `level' & \I{l} \\
  \B{m} & as in `man' or `from' & \I{m} \\
  \B{n} & as in `now' or `on' & \I{n} \\
  \B{ng} & one sound as in `singer', no audible g sound & \I{N} \\
  \B{p} & unaspirated as the p in `sport' at the start of a syllable & \I{p} \\
             & unreleased at the end of a syllable & \I{p\textcorner} \\
  \B{px} & pronounced p while holding breath & \I{p'} \\ 
  \B{r} & fronted, flapped r, as in Spanish `torero' & \I{R} \\
  \B{s} & sharp s sound as in `soft' or `floss' & \I{s} \\
  \B{t} & unaspirated as the t in `stop' at the start of a syllable & \I{t} \\
             & unreleased at the end of a syllable & \I{t\textcorner} \\
  \B{tx} & pronounced t while holding breath & \I{t'} \\
  \B{ts} & as the ts in `tsunami' or `cats' & \I{\t{ts}} \\
  \B{v} & as in `very' or `have' & \I{v} \\
  \B{w} & as in `what' or `will' & \I{w} \\
  \B{y} & as in `yet' & \I{j} \\ 
  \B{z} & as in `zoom' or `Oz' & \I{z} \\
\end{tabular} \\

\subsubsection*{Stress}

Word stress in Na'vi is unpredictable. It must be learned. Once learned it has the advantage that the stress usually remains on the same syllable no matter whether affixes are added. With mono\hyp{}syllabic words the primary stress is clear and has not been marked in the IPA. The stress in poly\hyp{}syllabic words is indicated by color marking in the main entry, e.g. \B{\cp{ta}ron}. For those comfortable with the IPA the stress mark ( \B{\I{"}} ) is used, e.g. \I{["ta.Ron]}.

\subsubsection*{Sound Change -- Lenition}

Lenition is a phonological change from one sound into another. Prefixes or adpositions that cause lenition are marked with a \texttt{+}. In Na'vi the following consonants undergo lenition: 
\begin{center}
\begin{tabular}{ l c l || l c l || l c l || l c l }
   \B{'} & \ding{234} & (\emph{disappears}) & \B{kx} & \ding{234} & \B{k} & \B{px} & \ding{234} & \B{p} & \B{tx} & \ding{234} &  \B{t} \\ 
   \B{k} & \ding{234} & \B{h} & \B{p} & \ding{234} & \B{f} & \B{t} & \ding{234} & \B{s} & \B{ts} & \ding{234} & \B{s} \\
\end{tabular} 
\end{center}

This is important to know since sometimes a word in use could be lenited and therefore is not found under the subsequent entry but rather under its unlenited rootform. Be aware of that.

\subsection*{A Note on Dialects}

Although Forest Na'vi is chosen for this dictionary as a point from which to start, it is by no means preferred or superior to any other dialect spoken on Pandora. The known dialects so far are mutually comprehensible, which means that it should not be difficult to switch between one dialect and another.

\subsubsection*{Of Forest Na'vi (Lì'fya Na'rìngä)}

This dialect is spoken by the Omatikaya Clan. So far, it is the best studied and well\hyp{}known variety of the Na'vi language and is therefore chosen to represent Na'vi in this dictionary as a point of reference.  

\subsubsection*{Of Reef Na'vi (Lì'fya Wionä)} 

This dialect is spoken by the Reef Clans like the Metkayina. It differs from Forest Na'vi in the following aspects:

\begin{itemize}

	\item The consonant cluster sy- \I{[sj]} is pronounced \I{[S]}, like the sh in `shore' or `ash'. The words \B{syawm}, \emph{know}, \B{syaw}, \emph{call}, or \B{syeha}, \emph{breath}, are pronounced \I{[Sawm]}, \I{[Saw]}, and \I{["SE.ha]} respectively in Reef Na'vi.
	\item The consonant cluster tsy- \I{[\t{ts}j]} is pronounced \I{[\t{tS}]}, like the ch in `church'. Thus the words \B{tsyal}, \emph{wing}, or \B{tsyeytsyìp}, \emph{tiny bite}, are pronounced \I{[\t{tS}al]} and \I{["\t{tS}Ej.\t{tS}Ip\textcorner]} respectively in Reef Na'vi.
	\item The glottal stop is lost beween non\hyp{}identical vowels, e.g. \B{fra'u}, \emph{everything}, becomes \B{frau}; \B{Lo'ak}, \emph{the name Lo'ak}, becomes \B{Loak}.
	\item The glottal stop between identical vowels \emph{may be} left out, e.g. \B{rä'ä}, \emph{do not}, may become \B{rää} \I{[R\ae."{}\ae]}; \B{me'em}, \emph{harmonious}, may become \B{meem} \I{[mE."{}Em]}.
		\subitem Given that speakers of Reef Na'vi are used to hearing these sequences of identical vowels, forms like \B{apxaa} \I{[a."ba.a]}, \B{meeveng} \I{[mE."{}E.vEN]}, \B{seii} \I{[sE.i."{}i]}, and \B{peekxinùm} \I{[pE.E."gi.nUm]} remain as they are. 
	\item At the beginning of a syllable, the ejectives \B{px} \I{[p']}, \B{tx} \I{[t']}, and \B{kx} \I{[k']} are pronounced \I{[b]}, \I{[d]}, and \I{[g]} respectively in this dialect, e.g. \B{txon} \I{[don]}, \emph{night}; \B{holpxay} \I{[hol."baj]}, \emph{number}, or \B{kxitx} \I{[git']}, \emph{death}.
		\subitem If the word ends in an ejective, the syllable boundaries will be reanalysed when a suffix is added, e.g. \B{'awkx} \I{[Pawk']}, but \B{'awkxìl} \I{["Paw.gIl]}
		\subitem When two ejectives across syllable boundaries get in contact within a word, they undergo regressive assimilation, e.g. \B{atxkxe} and \B{ekxtxu} are \I{[ad."gE]} and \I{[Eg."du]} rather than *\I{[at'."gE]} and *\I{[Ek'."du]}.
	\item Reef Na'vi distinguishes between a tense u sound \I{[u]} as in `food', and a lax u sound \I{[U]} as in `good'. Analogous to i and ì in Forest Na'vi, the sounds are marked u and ù, e.g. \B{pum} \I{[pum]} \emph{or} \I{[pUm]} (\textsc{rn:} \B{pùm} \I{[pUm]}) \emph{``dummy noun''}. This can also lead to minimal pairs in RN, e.g. \B{tsun} \I{[\t{ts}un]}\textsuperscript{\textsc{rn}} \emph{heel} and \B{tsùn} \I{[\t{ts}Un]}\textsuperscript{\textsc{rn}} \emph{can, be able}. These differences in sound are only marked in the respective entries for words with the lax ù sound. Assume that all other u sounds are tense in Reef Na'vi.
	\item In unstressed syllables \B{ä} \I{[\ae]} tends to change to \B{e} \I{[E]} in Reef Na'vi, e.g. \B{nge\-yä} \I{["NE.j\ae]} and \B{tä\-txaw} \I{[t\ae."t'aw]} may be pronounced \I{["NE.jE]} and \I{[tE."daw]}, respectively.
	\item The patientive ending \B{-ti} is favored over its other variants and the topical case \B{-(ì)ri} may also appear sentence final. 

\end{itemize}

% \B{} \I{[]} \emph{}

